% BXH-switches.tex -- testing successive switches with BXH and BXLE.
% Build: xelatex BXH-switches.tex (twice).
\documentclass[11pt]{article}
\usepackage{fontspec}
\usepackage{polyglossia}
\setdefaultlanguage{english}
\setotherlanguage{hebrew}
\setmainfont{Times New Roman}
\newfontfamily\hebrewfont[Script=Hebrew]{Times New Roman}
\setmonofont{Consolas}[Scale=MatchLowercase]
\usepackage[margin=1in]{geometry}
\usepackage{amsmath}
\usepackage{booktabs}
\usepackage{tikz}
\usetikzlibrary{arrows.meta,positioning,calc}
\usepackage{fancyvrb}
\usepackage[hidelinks]{hyperref}

\title{Testing Successive Switches with BXH and BXLE}
\author{Seymour J. Metz (H: Shmuel \texthebrew{שמואל בן לייביש ולאה})}
\date{10 November 2026}

% One register cell for the diagrams: \cell{x}{y}{text}{fill}
\newcommand{\cell}[4]{%
  \node[draw, minimum width=7mm, minimum height=7mm, fill=#4,
        font=\ttfamily\small] at (#1,#2) {#3};}

\begin{document}
\maketitle

\begin{abstract}
This paper explains William Collier's technique~\cite{collier1966} for
testing a string of switches with a sequence of BXH or BXLE instructions. It
shows why the technique works, the layout of the switch word it requires
(including a trailing guard bit that must be one), and its extension to
more modern processors.
\end{abstract}

\section{Origin}

The technique is due to William Collier, who described it in IBM Technical
Report TR 00.1412-1, dated 22 June 1966~\cite{collier1966}.

\section{The instructions}

While BXH and BXLE were designed for loop control, specifying R1 and R3 as
the same odd register allows their use for testing 30 consecutive switches in
succession.

\section{The trick}

With $R_1 = R_3$ odd, the instructions double the contents $x$, compare the
new contents $2x$ to the original contents $x$, and branch depending on the
compare result. Figure~\ref{fig:one} shows one execution.

\begin{figure}[ht]
\centering
\begin{tikzpicture}[x=7mm,y=7mm]
  \node[anchor=east, font=\small] at (-0.8,0) {$x$ (before)};
  \foreach \b [count=\i from 0] in {s,b,?,?,?,?,?,?} {
    \ifnum\i=0 \cell{\i}{0}{\b}{gray!25}
    \else\ifnum\i=1 \cell{\i}{0}{\b}{yellow!40}
    \else \cell{\i}{0}{\b}{white}\fi\fi }
  \node[anchor=east, font=\small] at (-0.8,-2.4) {$2x$ (after)};
  \foreach \b [count=\i from 0] in {b,?,?,?,?,?,?,?} {
    \ifnum\i=0 \cell{\i}{-2.4}{\b}{yellow!40}
    \else \cell{\i}{-2.4}{\b}{white}\fi }
  \foreach \i in {1,...,7} {
    \draw[-{Stealth}, thick] (\i,-0.45) -- (\i-1,-1.95); }
  \node[font=\small, align=left, anchor=west] at (8.3,-1.2)
    {every bit moves one place left;\\
     $s$ (the old sign) is lost;\\
     $b$ becomes the new sign;\\
     the comparand is still the old $x$};
\end{tikzpicture}
\caption{One execution of \texttt{BXH r,r,target}: the register is doubled,
that is, shifted left one bit (8 bits shown).}
\label{fig:one}
\end{figure}

\section{Why it tests one bit}

Let $s$ be the sign bit of $x$ and $b$ the bit that follows it, the bit that
becomes the sign of $2x$. Table~\ref{tab:cases} takes the four cases.

\begin{table}[ht]
\centering
\begin{tabular}{ccll}
\toprule
$s$ & $b$ & $2x$ (overflow ignored) & is $2x > x$? \\
\midrule
0 & 0 & non-negative, no overflow & yes \\
0 & 1 & negative (overflow) & no: $2x < 0 \le x$ \\
1 & 0 & non-negative (overflow) & yes: $2x \ge 0 > x$ \\
1 & 1 & negative, no overflow & no: $2x = x + x < x$ \\
\bottomrule
\end{tabular}
\caption{The four cases of one BXH execution.}
\label{tab:cases}
\end{table}

For nonzero $x$, BXH branches when $b = 0$ and BXLE branches when $b = 1$.
Because each
execution leaves the shifted value in the register, the next execution tests
the next bit. A sequence of BXH instructions therefore reads the switches
one at a time, from the high-order end.

\section{The switch word}

Two positions in the word are not switches (Figure~\ref{fig:layout}).
\begin{itemize}
\item \textbf{The leading bit} (bit~0) is never tested: the first execution
  tests bit~1. Its value does not matter.
\item \textbf{The trailing bit} (bit~31/63) must be \textbf{one}. It guarantees
  that the register is never zero while switches remain to be tested, which
  rules out the exception above. With a zero there, a switch that is zero
  and is followed only by zero switches would read as one.
\end{itemize}
That leaves bits 1 to 30 for thirty switches in a fullword register. The grande forms
(BXHG, BXLEG, JXHG and JXLEG) give 62 switches in a grande register;
Section~\ref{sec:code}
shows both widths.

\begin{figure}[ht]
\centering
\begin{tikzpicture}[x=7mm,y=7mm]
  \cell{0}{0}{0}{gray!25}
  \foreach \i in {1,...,6} { \cell{\i}{0}{$f_{\i}$}{yellow!25} }
  \node at (7.5,0) {$\cdots$};
  \cell{9}{0}{$f_{29}$}{yellow!25}
  \cell{10}{0}{$f_{30}$}{yellow!25}
  \cell{11}{0}{1}{green!25}
  \node[font=\scriptsize] at (0,0.75) {0};
  \node[font=\scriptsize] at (1,0.75) {1};
  \node[font=\scriptsize] at (10,0.75) {30};
  \node[font=\scriptsize] at (11,0.75) {31};
  \draw (0,-0.6) -- (0,-1.1) node[below, font=\small, align=center]
    {leading bit\\(not tested)};
  \draw (5.5,-0.6) -- (5.5,-1.1) node[below, font=\small] {30 switches};
  \draw (11,-0.6) -- (11,-1.1) node[below, font=\small, align=center]
    {trailing guard\\(must be 1)};
\end{tikzpicture}
\caption{The fullword of switches.}
\label{fig:layout}
\end{figure}

\section{A worked example}

Table~\ref{tab:walk} follows a register through six tests, shown at 8 bits
so that every bit is visible: a leading 0, the switches 1~0~1~1~0~0, and the
trailing guard~1.

\begin{table}[ht]
\centering
\begin{tabular}{clrlrcc}
\toprule
test & $x$ (bits) & $x$ & $2x$ (bits) & $2x$ & BXH & switch \\
\midrule
1 & \texttt{01011001} &   89 & \texttt{10110010} &  $-78$ & falls through & 1 \\
2 & \texttt{10110010} &  $-78$ & \texttt{01100100} &  100 & branches & 0 \\
3 & \texttt{01100100} &  100 & \texttt{11001000} &  $-56$ & falls through & 1 \\
4 & \texttt{11001000} &  $-56$ & \texttt{10010000} & $-112$ & falls through & 1 \\
5 & \texttt{10010000} & $-112$ & \texttt{00100000} &   32 & branches & 0 \\
6 & \texttt{00100000} &   32 & \texttt{01000000} &   64 & branches & 0 \\
\bottomrule
\end{tabular}
\caption{Six successive executions of \texttt{BXH r,r,target} (8-bit
illustration). Wherever the sign changes (tests 1, 2, 3 and 5) the addition
overflows, and the overflow is ignored. In tests 5 and 6 only the trailing
guard keeps the register non-zero.}
\label{tab:walk}
\end{table}

\section{In code}\label{sec:code}

Since BXH branches when the switch is zero, it can branch around the code
for a switch that is off. BXLE branches when the switch is on. Every test is
executed in turn, whatever the earlier results; the sequence must run in
order, because each test depends on the shift performed by the one before.

The IF macro of the High Level Assembler Toolkit's structured programming
macros has no syntax for BXH, BXHG, BXLE, BXLEG, JXH, JXHG, JXLE or JXLEG.

\subsection{32-bit instructions: 30 switches in a fullword register}

\begin{Verbatim}[fontsize=\small]
         L     R15,STATUS         odd register: R1 = R3 = R15
         JXH   R15,R15,SKIP1      switch 1 off: skip action 1
*        ... action 1 ...
SKIP1    JXH   R15,R15,SKIP2      switch 2 off: skip action 2
*        ... action 2 ...
SKIP2    DS    0H
*        ...
STATUS   DC    AL.1(0,FLAG1,...,FLAG30,1)
\end{Verbatim}

JXLE branches to the code for a switch that is on:

\begin{Verbatim}[fontsize=\small]
         JXLE  R15,R15,DO1        switch 1 on: go to action 1
\end{Verbatim}

\subsection{Grande instructions: 62 switches in a grande register}

The grande forms JXHG and JXLEG use the whole 64-bit register, so a grande
register holds a leading bit, 62 switches and the trailing guard.

\begin{Verbatim}[fontsize=\small]
         LG    R15,STATUS64       odd register: R1 = R3 = R15
         JXHG  R15,R15,SKIP1      switch 1 off: skip action 1
*        ... action 1 ...
SKIP1    JXHG  R15,R15,SKIP2      switch 2 off: skip action 2
*        ... action 2 ...
SKIP2    DS    0H
*        ...
         DS    0D
STATUS64 DC    AL.1(0,FLAG1,...,FLAG62,1)
\end{Verbatim}

JXLEG branches when the switch is on:

\begin{Verbatim}[fontsize=\small]
         JXLEG R15,R15,DO1        switch 1 on: go to action 1
\end{Verbatim}

\section*{Acknowledgments}
The technique and its first description are William Collier's. This paper
was drafted with the assistance of Claude (Anthropic), an AI engine.

\begin{thebibliography}{9}
\bibitem{collier1966}
William Collier, IBM Technical Report TR 00.1412-1, 22 June 1966.
\end{thebibliography}

\end{document}
